% =============================================================================
% Chapter 1 — Zero Page ($0000–$00FF)
% =============================================================================
\chapter{Zero Page}

\epigraph{``The zero page is prime real estate on the 6502. If your code isn't
living there, it's commuting.''}{\textit{--- Unattributed, ca.\ 1983}}

\section*{Introduction}

On the MOS 6502, the 256 bytes from \$0000 to \$00FF enjoy a privileged
status.  Instructions that reference zero-page addresses are one byte shorter
and one cycle faster than their absolute-address equivalents.  The CPU's
limited register set makes these locations indispensable --- they serve as
extra registers, scratch space, and indirect-address pointers for nearly every
significant operation.

In the NovaVM, the zero page is owned almost entirely by EhBASIC.  The
interpreter packs its warm-start vector, terminal state, memory-management
pointers, floating-point accumulators, and even a small machine-language
subroutine into this single page.  Understanding these locations is the key to
understanding how BASIC works under the hood --- and to writing programs that
can peek behind the curtain.

The map below is arranged in address order.  Where a location serves double
duty under different aliases (a common EhBASIC technique), all names are
listed.

% ═══════════════════════════════════════════════════════════════════════════════
\section{Boot and Warm Start (\$00--\$02)}
% ═══════════════════════════════════════════════════════════════════════════════

When EhBASIC starts up, it writes a three-byte \texttt{JMP} instruction at the
very bottom of the zero page.  Thereafter, a warm start (the BASIC command
\texttt{RUN} completing, or the user pressing the reset key) simply jumps to
address \$0000, which executes the \texttt{JMP} and re-enters the interpreter's
main loop without losing the stored program.

\mementry{0000}{LAB\_WARM}{R/W}{\$4C}

The first byte of the warm-start trampoline.  Always contains \texttt{\$4C},
the 6502 opcode for \texttt{JMP abs}.  Do not alter this byte unless you are
deliberately redirecting the warm-start path.

\mementry{0001}{Wrmjpl}{R/W}{---}

Low byte of the warm-start jump target address.

\mementry{0002}{Wrmjph}{R/W}{---}

High byte of the warm-start jump target address.  Together with \addr{0001},
this forms the destination of the \texttt{JMP} at \addr{0000}.

\begin{notebox}
You can redirect the warm start by POKEing a new address into \$01--\$02.
For example, to jump to your own machine-language routine at \$8000:
\begin{lstlisting}
DOKE $01,$8000
\end{lstlisting}
Be very sure you know what you are doing --- a bad vector here means a
hard reset is the only way back.
\end{notebox}

% ═══════════════════════════════════════════════════════════════════════════════
\section{Unallocated (\$03--\$09)}
% ═══════════════════════════════════════════════════════════════════════════════

Addresses \$03 through \$09 are not assigned by EhBASIC.  They sit in the gap
between the warm-start trampoline and the USR() vector.  These seven bytes are
safe for your own machine-language routines or short data tables, provided you
do not call any BASIC routine that might overwrite them during
initialisation.  In practice, they remain untouched after boot.

% ═══════════════════════════════════════════════════════════════════════════════
\section{USR() Function Vector (\$0A--\$0C)}
% ═══════════════════════════════════════════════════════════════════════════════

The \texttt{USR()} function in EhBASIC calls a user-supplied machine-language
routine.  Like the warm-start vector, the mechanism is a three-byte
\texttt{JMP} instruction planted in the zero page.

\mementry{000A}{Usrjmp}{R/W}{\$4C}

The \texttt{JMP} opcode byte for the USR() dispatch.  Always \$4C.

\mementry{000B}{Usrjpl}{R/W}{---}

Low byte of the USR() target address.

\mementry{000C}{Usrjph}{R/W}{---}

High byte of the USR() target address.  To point USR() at a routine at
\$4000:

\begin{lstlisting}
DOKE $0B,$4000
X = USR(42)
\end{lstlisting}

The argument (42 in this case) is passed in FAC1 (see \addr{00AC}), and the
routine should leave its return value there.

% ═══════════════════════════════════════════════════════════════════════════════
\section{Terminal Settings (\$0D--\$10)}
% ═══════════════════════════════════════════════════════════════════════════════

These four bytes control how EhBASIC interacts with the terminal --- line
width, cursor column tracking, and null padding.

\mementry{000D}{Nullct}{R/W}{\$00}

Null count: the number of \$00 bytes transmitted after each carriage return.
On real serial terminals this gave the print head time to return to the left
margin.  On the NovaVM this is unused and should be left at zero.

\mementry{000E}{TPos}{R/W}{\$00}

Terminal position: the current column of the cursor, maintained by EhBASIC's
output routines.  Starts at zero for the left margin and increments with each
printed character.  The interpreter checks this against \addr{000F} to decide
when to wrap.  Reading this byte tells you the cursor's horizontal position:

\begin{lstlisting}
PRINT "Hello";
PRINT PEEK($0E)
\end{lstlisting}

This prints \texttt{5}, the column after the five characters of ``Hello''.

\mementry{000F}{TWidth}{R/W}{\$50}

Terminal width.  Set to 80 (\$50) by the NovaVM at boot.  EhBASIC wraps
output and inserts a carriage return when \addr{000E} reaches this value.

\begin{lstlisting}
POKE $0F,40 : REM narrow screen
\end{lstlisting}

\begin{warningbox}
Changing TWidth does not resize the physical display.  It only affects
EhBASIC's line-wrap logic.  Setting it wider than 80 will cause output to
overrun the visible screen.
\end{warningbox}

\mementry{0010}{Iclim}{R/W}{\$50}

Input column limit: the maximum number of characters EhBASIC will accept on a
single \texttt{INPUT} line.  Normally matches TWidth.

% ═══════════════════════════════════════════════════════════════════════════════
\section{Temporary Integer (\$11--\$12)}
% ═══════════════════════════════════════════════════════════════════════════════

\mementry{0011}{Itempl}{R/W}{---}

Low byte of a temporary 16-bit integer used internally during line-number
parsing, number-base conversion, and various arithmetic operations.  Also
known as \texttt{nums\_1} when used for binary/hex string conversion.

\mementry{0012}{Itemph}{R/W}{---}

High byte of the temporary integer.  Also \texttt{nums\_2}.  A third alias
byte \texttt{nums\_3} occupies \$13 for 24-bit conversions.

These locations are transient workspace.  Their contents are unpredictable
between BASIC statements and should not be relied upon for storage.

% ═══════════════════════════════════════════════════════════════════════════════
\section{Unallocated (\$13--\$5A)}
% ═══════════════════════════════════════════════════════════════════════════════

Addresses \$13 through \$5A --- a full 72 bytes --- are not assigned by
EhBASIC.  This is the largest contiguous free region on the zero page and an
excellent place to store variables, counters, or short machine-language
routines that benefit from fast zero-page addressing.

\begin{tipbox}
If you are writing a machine-language game loop or interrupt handler, the
\$13--\$5A block gives you plenty of room for sprite coordinates, frame
counters, velocity tables, and other frequently-accessed state --- all at
zero-page speed.
\end{tipbox}

% ═══════════════════════════════════════════════════════════════════════════════
\section{Interpreter Flags (\$5B--\$64)}
% ═══════════════════════════════════════════════════════════════════════════════

This cluster of bytes controls the interpreter's parsing and evaluation state.
Most are written and read only by internal BASIC routines; they are documented
here for completeness and for advanced users writing machine-language
extensions.

\mementry{005B}{Srchc / Temp3}{R/W}{---}

Search character used by the tokeniser and line-editing routines.  Also serves
as \texttt{Temp3}, a scratch byte during number-to-string conversion.  Double
duty as the low byte of \texttt{XOAw} (the EOR/OR/AND working word).

\mementry{005C}{Scnquo / Asrch}{R/W}{---}

Scan-between-quotes flag.  When set, the tokeniser treats everything up to the
next quotation mark as literal text rather than keywords.  Also used as
\texttt{Asrch} (alternate search character) and as the high byte of
\texttt{XOAw}.

\mementry{005D}{Ibptr / Dimcnt / Tindx}{R/W}{---}

A three-way alias:
\begin{itemize}
  \item \textbf{Ibptr} --- pointer (offset) into the input buffer during line entry.
  \item \textbf{Dimcnt} --- dimension count while processing \texttt{DIM} statements.
  \item \textbf{Tindx} --- token index during tokenisation.
\end{itemize}
Only one use is active at a time; the interpreter knows which context applies.

\mementry{005E}{Defdim}{R/W}{\$00}

Default DIM flag.  When non-zero, an array reference that has not been
explicitly dimensioned will be auto-dimensioned (to a default upper bound of
10).  Set internally by the array-access routines.

\mementry{005F}{Dtypef}{R/W}{\$00}

Data type flag.  \$FF means the current value is a string; \$00 means numeric.
The interpreter sets this after every expression evaluation so that subsequent
operations (assignment, PRINT, comparison) know which type they are handling.

\mementry{0060}{Oquote / Gclctd}{R/W}{\$00}

Open-quote flag during DATA scanning and LIST output.  Bit 7 signals that the
scanner is inside a quoted string.  The same byte doubles as \texttt{Gclctd},
the garbage-collected flag --- set after the string heap has been compacted to
prevent redundant re-collection.

\mementry{0061}{Sufnxf}{R/W}{\$00}

Subscript/FN flag.  Bit 7 set indicates a \texttt{FN} call is in progress;
the lower bits track subscript depth.  This prevents the interpreter from
confusing array subscripts with function arguments.

\mementry{0062}{Imode}{R/W}{\$00}

Input mode flag.  \$00 during an \texttt{INPUT} statement, \$80 during
\texttt{READ}.  This tells the shared input-parsing code whether to prompt the
user or consume the next \texttt{DATA} element.

\mementry{0063}{Cflag}{R/W}{\$00}

Comparison evaluation flag.  Accumulates the result of relational operators
during expression parsing:
\begin{itemize}
  \item Bit 0 --- less than (\texttt{<})
  \item Bit 1 --- equal (\texttt{=})
  \item Bit 2 --- greater than (\texttt{>})
\end{itemize}
Combinations like \texttt{<=} set bits 0 and 1 simultaneously.

\mementry{0064}{TabSiz}{R/W}{\$0A}

TAB step size.  Controls the column spacing of comma-separated items in
\texttt{PRINT} statements.  Default is 10.

\begin{lstlisting}
POKE $64,20 : REM wider tab stops
PRINT 1,2,3
\end{lstlisting}

% ═══════════════════════════════════════════════════════════════════════════════
\section{String Descriptor Stack (\$65--\$70)}
% ═══════════════════════════════════════════════════════════════════════════════

EhBASIC maintains a small stack of string descriptors in the zero page.  Each
descriptor is three bytes: one byte for the string length and two bytes for
the pointer to the string data.  The stack holds up to three temporary strings
at once.

\mementry{0065}{next\_s}{R/W}{---}

Pointer to the next available slot on the descriptor stack.  Starts at
\$68 (the base of \texttt{des\_sk}) and advances by three for each temporary
string pushed.

\mementry{0066}{last\_sl}{R/W}{---}

Low byte of the address of the most recently pushed descriptor.

\mementry{0067}{last\_sh}{R/W}{\$00}

High byte of the last descriptor address.  Since the descriptor stack lives
entirely within the zero page, this byte is always \$00.

\memrange{0068}{0070}{des\_sk (Descriptor Stack)}{R/W}{---}

Nine bytes holding up to three string descriptors.  These are working storage
for intermediate string results --- for example, the expression
\texttt{A\$+B\$+C\$} will push descriptors onto this stack as it concatenates.
If more than three temporaries are needed, EhBASIC raises a
\texttt{?FORMULA TOO COMPLEX} error.

% ═══════════════════════════════════════════════════════════════════════════════
\section{Utility Pointers and FAC Temporaries (\$71--\$78)}
% ═══════════════════════════════════════════════════════════════════════════════

\mementry{0071}{ut1\_pl}{R/W}{---}

Utility pointer 1, low byte.  A general-purpose 16-bit pointer used by block
move operations, array manipulation, and string handling.  Also known as
\texttt{Temp\_2} during block moves.

\mementry{0072}{ut1\_ph}{R/W}{---}

Utility pointer 1, high byte.

\mementry{0073}{ut2\_pl}{R/W}{---}

Utility pointer 2, low byte.  A second general-purpose pointer, often used as
the destination address when \texttt{ut1} is the source.

\mementry{0074}{ut2\_ph}{R/W}{---}

Utility pointer 2, high byte.

\mementry{0075}{FACt\_1}{R/W}{---}

FAC temporary mantissa byte 1.  Used as scratch space during floating-point
operations and multiply/divide routines.

\mementry{0076}{FACt\_2 / dims\_l}{R/W}{---}

FAC temporary mantissa byte 2.  Also serves as \texttt{dims\_l}, the low byte
of the array dimension size during \texttt{DIM} processing.

\mementry{0077}{FACt\_3 / dims\_h}{R/W}{---}

FAC temporary mantissa byte 3, and the high byte of the array dimension size.

\mementry{0078}{TempB}{R/W}{---}

A single scratch byte used by various interpreter routines.  Entirely
transient.

% ═══════════════════════════════════════════════════════════════════════════════
\section{Memory Management Pointers (\$79--\$86)}
% ═══════════════════════════════════════════════════════════════════════════════

These seven pointer pairs are the backbone of EhBASIC's memory management.
They define the boundaries of every major region in RAM: program text,
variables, arrays, and string storage.  Understanding them is essential for
anyone who wants to know how much memory is free, where a program lives, or
how to relocate data.

The pointers are arranged so that, in a healthy system:

\begin{center}
\texttt{Smem} $\leq$ \texttt{Svar} $\leq$ \texttt{Sarry} $\leq$
\texttt{Earry} $\leq$ \texttt{Sstor} $\leq$ \texttt{Emem}
\end{center}

Program text grows upward from Smem; string storage grows \emph{downward} from
Emem.  The gap between Earry and Sstor is free memory.

\mementry{0079}{Smeml}{R/W}{---}

Start of BASIC program text, low byte.

\mementry{007A}{Smemh}{R/W}{---}

Start of BASIC program text, high byte.  This is the base address of the
tokenised program.  On the NovaVM, this is typically \$0280 (the first byte
above the vector table).

\begin{lstlisting}
PRINT DEEK($79)
\end{lstlisting}

This prints the start-of-program address in decimal.

\mementry{007B}{Svarl}{R/W}{---}

Start of variables, low byte.

\mementry{007C}{Svarh}{R/W}{---}

Start of variables, high byte.  This pointer also marks the end of program
text.  The difference \texttt{DEEK(\$7B) - DEEK(\$79)} is the size of your
program in bytes.

\begin{lstlisting}
PRINT DEEK($7B)-DEEK($79);"BYTES PROGRAM"
\end{lstlisting}

\mementry{007D}{Sarryl}{R/W}{---}

Start of arrays, low byte.

\mementry{007E}{Sarryh}{R/W}{---}

Start of arrays, high byte.  Simple variables (scalars) occupy the region
from Svar to Sarry.

\mementry{007F}{Earryl}{R/W}{---}

End of arrays, low byte.

\mementry{0080}{Earryh}{R/W}{---}

End of arrays, high byte.  All dimensioned arrays live between Sarry and
Earry.

\mementry{0081}{Sstorl}{R/W}{---}

Bottom of string storage, low byte.

\mementry{0082}{Sstorh}{R/W}{---}

Bottom of string storage, high byte.  String data is allocated from the top of
memory downward.  Sstor marks the lowest string address currently in use.  As
more strings are created, this pointer moves down toward Earry.  When the two
meet, you get an \texttt{?OUT OF MEMORY} error.

\mementry{0083}{Sutill}{R/W}{---}

String utility pointer, low byte.

\mementry{0084}{Sutilh}{R/W}{---}

String utility pointer, high byte.  A working pointer used during string
allocation and garbage collection.

\mementry{0085}{Ememl}{R/W}{---}

End of memory, low byte.

\mementry{0086}{Ememh}{R/W}{---}

End of memory, high byte.  This is the top of available RAM --- the highest
address BASIC is allowed to use.  String storage begins just below this point
and grows downward.

\begin{lstlisting}
REM -- Free memory calculation --
FM = DEEK($81)-DEEK($7F)
PRINT FM;"BYTES FREE"
\end{lstlisting}

\begin{notebox}
The FRE() function in EhBASIC returns the same value but also triggers a
garbage collection pass, which can cause a brief pause.  Reading the pointers
directly with DEEK is instantaneous but does not account for reclaimable
string space.
\end{notebox}

\begin{tipbox}
A quick memory map dump:
\begin{lstlisting}
PRINT "PROGRAM  ";DEEK($79);" TO ";DEEK($7B)-1
PRINT "VARIABLES";DEEK($7B);" TO ";DEEK($7D)-1
PRINT "ARRAYS   ";DEEK($7D);" TO ";DEEK($7F)-1
PRINT "FREE     ";DEEK($7F);" TO ";DEEK($81)-1
PRINT "STRINGS  ";DEEK($81);" TO ";DEEK($85)-1
\end{lstlisting}
\end{tipbox}

% ═══════════════════════════════════════════════════════════════════════════════
\section{Line and Data Pointers (\$87--\$9F)}
% ═══════════════════════════════════════════════════════════════════════════════

These pointers track the interpreter's current position in the program, the
DATA read cursor, and the state of the most recent variable reference.

\mementry{0087}{Clinel}{R/W}{---}

Current BASIC line number, low byte.

\mementry{0088}{Clineh}{R/W}{---}

Current BASIC line number, high byte.  While a program is running, this pair
holds the line number currently being executed.  In immediate mode, it
contains \$FFFF.

\begin{lstlisting}
10 PRINT "LINE ";DEEK($87)
20 PRINT "LINE ";DEEK($87)
RUN
\end{lstlisting}

This prints \texttt{LINE 10} then \texttt{LINE 20}.

\mementry{0089}{Blinel}{R/W}{---}

Break line number, low byte.

\mementry{008A}{Blineh}{R/W}{---}

Break line number, high byte.  When a program is stopped (by the \texttt{STOP}
statement, Ctrl+C, or an error), EhBASIC saves the line number here.  The
\texttt{CONT} command uses this to resume execution.

\mementry{008B}{Cpntrl}{R/W}{---}

Continue pointer, low byte.

\mementry{008C}{Cpntrh}{R/W}{---}

Continue pointer, high byte.  The address within program text where execution
should resume after \texttt{CONT}.  If you edit any line, this pointer is
invalidated and \texttt{CONT} will refuse to proceed.

\mementry{008D}{Dlinel}{R/W}{---}

Current DATA line number, low byte.

\mementry{008E}{Dlineh}{R/W}{---}

Current DATA line number, high byte.  Tracks which line the DATA cursor is on,
so that an \texttt{?OUT OF DATA} error can report the correct line number.

\mementry{008F}{Dptrl}{R/W}{---}

DATA pointer, low byte.

\mementry{0090}{Dptrh}{R/W}{---}

DATA pointer, high byte.  Points to the next byte to be consumed by a
\texttt{READ} statement.  \texttt{RESTORE} resets this to the start of program
text.

\mementry{0091}{Rdptrl}{R/W}{---}

Read pointer, low byte.

\mementry{0092}{Rdptrh}{R/W}{---}

Read pointer, high byte.  Used during INPUT/READ to track the current position
in the data being parsed.

\mementry{0093}{Varnm1}{R/W}{---}

Current variable name, first byte.  When the interpreter encounters a variable
reference (e.g.\ \texttt{A\$}), the first character of the name is stored here.

\mementry{0094}{Varnm2}{R/W}{---}

Current variable name, second byte.  For single-letter variables this is zero.
For two-letter names, it holds the second character.  Bit 7 of either byte may
be set to indicate a string (\$) or integer (\%) type suffix.

\mementry{0095}{Cvaral}{R/W}{---}

Current variable address, low byte.

\mementry{0096}{Cvarah}{R/W}{---}

Current variable address, high byte.  After a variable lookup, this pair
points to the variable's storage in the variable table (between Svar and
Sarry).

\mementry{0097}{Frnxtl}{R/W}{---}

FOR/NEXT variable pointer, low byte.  Also aliased as \texttt{Tidx1} (temp
line index) and \texttt{Lvarpl} (LET variable pointer).

\mementry{0098}{Frnxth}{R/W}{---}

FOR/NEXT variable pointer, high byte.  Also \texttt{Lvarph}.  During FOR/NEXT
execution, this pair points to the loop control variable so that NEXT can
find the matching FOR frame on the stack.

\mementry{0099}{prstk}{R/W}{\$00}

Precedence-stacked flag.  Set non-zero when the expression evaluator has
pushed an operator onto the precedence stack.  Cleared when the stack is
unwound.

\mementry{009B}{comp\_f}{R/W}{\$00}

Compare function flag.  Uses bits 0, 1, and 2 to encode the comparison
being evaluated:

\bitfield{---}{---}{---}{---}{---}{\textgreater}{=}{\textless}

This is the internal companion to \addr{0063} (Cflag), used during the
parsing phase of a comparison expression.

\mementry{009C}{func\_l / garb\_l}{R/W}{---}

Function pointer low byte.  During function evaluation, points to the body of
a \texttt{DEF FN} definition.  Also serves as \texttt{garb\_l}, the garbage
collection working pointer low byte.

\mementry{009D}{func\_h / garb\_h}{R/W}{---}

Function pointer high byte / garbage collection working pointer high byte.

\mementry{009E}{des\_2l}{R/W}{---}

String descriptor 2 pointer, low byte.

\mementry{009F}{des\_2h}{R/W}{---}

String descriptor 2 pointer, high byte.  Used during string operations that
need to reference two descriptors simultaneously (e.g.\ concatenation or
comparison).

% ═══════════════════════════════════════════════════════════════════════════════
\section{Garbage Collection and Array Workspace (\$A0--\$AB)}
% ═══════════════════════════════════════════════════════════════════════════════

\mementry{00A0}{g\_step}{R/W}{---}

Garbage collection step size.  Controls the search increment during string
heap compaction.

\mementry{00A1}{Fnxjmp}{R/W}{\$4C}

Jump vector for function dispatch --- another three-byte \texttt{JMP}
instruction, like the warm-start and USR() vectors.  Byte \$A1 holds the
\texttt{JMP} opcode.

\mementry{00A2}{Fnxjpl / g\_indx}{R/W}{---}

Low byte of the function jump target.  Also aliased as \texttt{g\_indx}, a
temporary index used during garbage collection.

\mementry{00A3}{FAC2\_r}{R/W}{---}

FAC2 rounding byte.  The extra precision byte for the second floating-point
accumulator, used during multi-step arithmetic to reduce rounding errors.

\mementry{00A4}{Adatal / Nbendl}{R/W}{---}

Array data pointer, low byte.  Also serves as \texttt{Nbendl} --- the new
block end pointer during memory block moves (INSERT, DELETE).

\mementry{00A5}{Adatah / Nbendh}{R/W}{---}

Array data pointer, high byte / new block end pointer high byte.

\mementry{00A6}{Obendl}{R/W}{---}

Old block end pointer, low byte.  Used alongside \texttt{Nbend} during block
move operations.

\mementry{00A7}{Obendh}{R/W}{---}

Old block end pointer, high byte.

\mementry{00A8}{numexp / numbit}{R/W}{---}

Number exponent count during string-to-float conversion.  Also \texttt{numbit},
a bit counter during array element address calculations.

\mementry{00A9}{expcnt}{R/W}{---}

Exponent count for string-to-float conversion --- tracks the decimal exponent
independently from the mantissa being assembled.

\mementry{00AA}{numdpf / Astrtl}{R/W}{---}

Decimal point flag during number conversion.  This byte is heavily aliased:
also \texttt{Astrtl} (array start pointer low), \texttt{Histrl} (highest
string low), \texttt{Baslnl} (BASIC search line pointer low),
\texttt{Fvar\_l} (find variable pointer low), \texttt{Ostrtl} (old block
start pointer low), and \texttt{Vrschl} (variable search pointer low).

Only one use is active at any given time --- the interpreter context determines
the meaning.

\mementry{00AB}{expneg / Astrth}{R/W}{---}

Exponent negate flag during number conversion.  Similarly aliased:
\texttt{Astrth}, \texttt{Histrh}, \texttt{Baslnh}, \texttt{Fvar\_h},
\texttt{Ostrth}, \texttt{Vrschh}.

% ═══════════════════════════════════════════════════════════════════════════════
\section{Floating-Point Accumulator 1 (\$AC--\$B2)}
% ═══════════════════════════════════════════════════════════════════════════════

FAC1 is the primary floating-point accumulator.  All arithmetic operations
ultimately produce their result here.  EhBASIC uses a five-byte format: one
byte for the exponent and four bytes for the mantissa (with an implicit
leading 1 bit) plus a separate sign byte.

\mementry{00AC}{FAC1\_e}{R/W}{---}

FAC1 exponent.  Uses excess-128 notation: the stored value minus 128 gives the
true binary exponent.  A stored zero means the entire FAC is zero.  Also
aliased as \texttt{str\_ln} (string length) during string operations.

\mementry{00AD}{FAC1\_1}{R/W}{---}

FAC1 mantissa byte 1 (most significant).  Also \texttt{str\_pl} (string
pointer low) in string context.

\mementry{00AE}{FAC1\_2}{R/W}{---}

FAC1 mantissa byte 2.  Also \texttt{str\_ph} (string pointer high) and
\texttt{des\_pl} (descriptor pointer low).

\mementry{00AF}{FAC1\_3}{R/W}{---}

FAC1 mantissa byte 3 (least significant in the four-byte mantissa).  Also
\texttt{des\_ph} (descriptor pointer high) and \texttt{mids\_l} (MID\$
temporary length).

\mementry{00B0}{FAC1\_s}{R/W}{---}

FAC1 sign.  Bit 7 indicates negative.  The remaining bits are unused.

\mementry{00B1}{negnum / numcon}{R/W}{---}

String-to-float negate flag: set when parsing a negative numeric literal.
Also \texttt{numcon}, the constant counter during polynomial series evaluation
(used by SIN, COS, EXP, LOG, ATN).

\mementry{00B2}{FAC1\_o}{R/W}{---}

FAC1 overflow byte.  Holds the extra precision bit that falls off the bottom
of the mantissa during normalisation.  Used for rounding the final result.

% ═══════════════════════════════════════════════════════════════════════════════
\section{Floating-Point Accumulator 2 (\$B3--\$B9)}
% ═══════════════════════════════════════════════════════════════════════════════

FAC2 is the secondary accumulator.  Binary operations (add, subtract, multiply,
divide) load one operand into FAC2 and the other into FAC1; the result lands
in FAC1.

\mementry{00B3}{FAC2\_e}{R/W}{---}

FAC2 exponent.  Same excess-128 encoding as FAC1.

\mementry{00B4}{FAC2\_1}{R/W}{---}

FAC2 mantissa byte 1 (most significant).

\mementry{00B5}{FAC2\_2}{R/W}{---}

FAC2 mantissa byte 2.

\mementry{00B6}{FAC2\_3}{R/W}{---}

FAC2 mantissa byte 3 (least significant).

\mementry{00B7}{FAC2\_s}{R/W}{---}

FAC2 sign byte.  Bit 7 indicates negative.

\mementry{00B8}{FAC\_sc}{R/W}{---}

FAC sign comparison result.  After comparing FAC1 and FAC2, this byte records
whether the signs agree or differ.  Also aliased as \texttt{ssptr\_l} (string
start pointer low) and \texttt{sdescr} (string descriptor pointer).

\mementry{00B9}{FAC1\_r}{R/W}{---}

FAC1 rounding byte.  The ninth bit of precision, used during final rounding of
arithmetic results.  Also \texttt{ssptr\_h} (string start pointer high).

% ═══════════════════════════════════════════════════════════════════════════════
\section{BASIC Byte Fetch Subroutines (\$BA--\$D7)}
% ═══════════════════════════════════════════════════════════════════════════════

This region contains one of EhBASIC's most distinctive features: executable
machine code planted directly in the zero page.  Two small subroutines live
here, giving the interpreter fast, one-byte \texttt{JSR} calls for its two
most frequent operations --- fetching the next byte from program text and
re-fetching the current byte.

\mementry{00BA}{csidx / Asptl}{R/W}{---}

Line crunch save index.  Used during tokenisation to remember the current
position.  Also aliased as \texttt{Asptl} (array size pointer low),
\texttt{Btmpl}, \texttt{Cptrl}, and \texttt{Sendl} --- all temporary pointer
low bytes used in different contexts.

\mementry{00BB}{Aspth}{R/W}{---}

Array size pointer high byte.  Also \texttt{Btmph}, \texttt{Cptrh}, and
\texttt{Sendh}.

\memrange{00BC}{00C1}{LAB\_IGBY (Increment and Get BYte)}{R/W}{---}

Six bytes of machine code implementing the ``get next BASIC byte'' subroutine.
When the interpreter calls \texttt{JSR \$BC}, this routine increments the
BASIC execute pointer (\addr{00C3}--\addr{00C4}) and falls through into
\texttt{LAB\_GBYT} to fetch the byte at the new address.

This subroutine is placed in the zero page purely for speed.  The 6502's
\texttt{JSR zp} is not actually faster (there is no zero-page JSR), but
having the code here keeps the execute pointer in the zero page where
\texttt{LDA (zp),Y} can reach it in just five cycles.

\memrange{00C2}{00C4}{LAB\_GBYT (Get BYTe)}{R/W}{---}

Three bytes of machine code that load the byte at the current BASIC execute
pointer and set the processor flags accordingly.  The interpreter calls
\texttt{JSR \$C2} whenever it needs to re-examine the current token without
advancing.

\mementry{00C3}{Bpntrl}{R/W}{---}

BASIC execute pointer, low byte.  This is the interpreter's program counter
--- it points to the byte in program text that will be (or has just been)
processed.

\mementry{00C4}{Bpntrh}{R/W}{---}

BASIC execute pointer, high byte.  Together with \addr{00C3}, this is the
single most important pointer in the entire interpreter.

\begin{notebox}
The region \$C5--\$D7 continues the LAB\_GBYT subroutine code.  These 19
bytes are part of the bulk-initialised machine code block and should never be
modified.  They do not correspond to individually meaningful variables.
\end{notebox}

% ═══════════════════════════════════════════════════════════════════════════════
\section{Pseudo-Random Number Generator (\$D8--\$DB)}
% ═══════════════════════════════════════════════════════════════════════════════

EhBASIC's \texttt{RND()} function uses a simple four-byte linear feedback
shift register stored in the zero page.

\mementry{00D8}{Rbyte4}{R/W}{---}

PRNG byte 4 (most significant).  The ``extra'' byte added by EhBASIC's
extended PRNG algorithm.

\mementry{00D9}{Rbyte1}{R/W}{---}

PRNG byte 1.

\mementry{00DA}{Rbyte2}{R/W}{---}

PRNG byte 2.

\mementry{00DB}{Rbyte3}{R/W}{---}

PRNG byte 3 (least significant).  Each call to \texttt{RND(1)} shifts and
XORs all four bytes to produce a new pseudo-random value.

\begin{lstlisting}
REM Seed the PRNG
X = RND(-42)
REM Now RND(1) gives a repeatable sequence
FOR I=1 TO 5 : PRINT RND(1) : NEXT
\end{lstlisting}

Calling \texttt{RND()} with a negative argument seeds the generator.
\texttt{RND(0)} returns the most recent value.  \texttt{RND(1)} advances
to the next.

% ═══════════════════════════════════════════════════════════════════════════════
\section{Interrupt Handler Vectors (\$DC--\$E1)}
% ═══════════════════════════════════════════════════════════════════════════════

The NovaVM extends EhBASIC with support for NMI and IRQ interrupts that can
be handled from BASIC.  Each handler has a control byte and a two-byte
address vector.

\mementry{00DC}{NmiBase}{R/W}{\$00}

NMI (Non-Maskable Interrupt) control register.

\bitfield{Enable}{Setup}{Fired}{---}{---}{---}{---}{---}

\begin{itemize}
  \item \textbf{Bit 7} --- Interrupt enabled.  When set, the NMI handler is armed.
  \item \textbf{Bit 6} --- Setup complete.  Set by the \texttt{NMI} BASIC command
        after installing a handler.
  \item \textbf{Bit 5} --- Interrupt occurred.  Set by the hardware when an NMI
        fires; cleared after the handler runs.
\end{itemize}

\mementry{00DD}{NMI Address Low}{R/W}{---}

Low byte of the NMI handler routine address.

\mementry{00DE}{NMI Address High}{R/W}{---}

High byte of the NMI handler address.  Together with \addr{00DD}, this points
to the BASIC line (or machine-language routine) that should execute when an
NMI arrives.

\mementry{00DF}{IrqBase}{R/W}{\$00}

IRQ (Interrupt Request) control register.  Same bit layout as \addr{00DC}.

\bitfield{Enable}{Setup}{Fired}{---}{---}{---}{---}{---}

\mementry{00E0}{IRQ Address Low}{R/W}{---}

Low byte of the IRQ handler routine address.

\mementry{00E1}{IRQ Address High}{R/W}{---}

High byte of the IRQ handler address.

\begin{lstlisting}
10 IRQ 1000
20 REM main loop here
30 GOTO 20
1000 REM interrupt handler
1010 PRINT "TICK"
1020 RETIRQ
\end{lstlisting}

The \texttt{IRQ} command sets up \addr{00DF}--\addr{00E1} automatically.  You
only need to POKE these directly if you are installing a machine-language
handler.

% ═══════════════════════════════════════════════════════════════════════════════
\section{HELP Scratch (\$E2)}
% ═══════════════════════════════════════════════════════════════════════════════

\mementry{00E2}{help\_len}{R/W}{\$00}

Scratch byte used by the \texttt{HELP} command to track the length of the
keyword being searched.  No user-visible purpose outside of that command.

% ═══════════════════════════════════════════════════════════════════════════════
\section{Unused (\$E3--\$EE)}
% ═══════════════════════════════════════════════════════════════════════════════

Addresses \$E3 through \$EE --- twelve bytes --- are completely unassigned.
They are safe for user programs, machine-language scratch space, or small
lookup tables.  Combined with the seven bytes at \$03--\$09 and the large
block at \$13--\$5A, this gives a total of 91 free zero-page bytes available
for your own use.

% ═══════════════════════════════════════════════════════════════════════════════
\section{Decimal Conversion Buffer (\$EF--\$FF)}
% ═══════════════════════════════════════════════════════════════════════════════

\memrange{00EF}{00FF}{Decss (Decimal String Buffer)}{R/W}{---}

Seventeen bytes reserved for number-to-decimal-string conversion.  When BASIC
needs to print a number or convert it with \texttt{STR\$()}, the result is
assembled here, character by character, starting from the least significant
digit and working left.  The buffer is large enough for the longest possible
numeric string: a sign, up to ten digits, a decimal point, an ``E'', an
exponent sign, and two exponent digits.

This buffer is transient --- its contents are only valid during and
immediately after a conversion.  Do not store persistent data here.

% ═══════════════════════════════════════════════════════════════════════════════
\section*{Summary}
% ═══════════════════════════════════════════════════════════════════════════════

The zero page is a tightly-packed workspace shared between EhBASIC and the
6502 hardware.  The table below summarises the allocation at a glance:

\begin{center}
\small
\begin{tabular}{@{}lll@{}}
\toprule
\textbf{Range} & \textbf{Usage} & \textbf{Owner} \\
\midrule
\$00--\$02     & Warm-start JMP               & EhBASIC \\
\$03--\$09     & \textit{Free}                & User \\
\$0A--\$0C     & USR() JMP vector             & EhBASIC \\
\$0D--\$12     & Terminal settings \& temp int & EhBASIC \\
\$13--\$5A     & \textit{Free (72 bytes)}     & User \\
\$5B--\$64     & Interpreter flags            & EhBASIC \\
\$65--\$70     & String descriptor stack       & EhBASIC \\
\$71--\$78     & Utility pointers \& FAC temp  & EhBASIC \\
\$79--\$86     & Memory management pointers   & EhBASIC \\
\$87--\$9F     & Line, data, \& var pointers  & EhBASIC \\
\$A0--\$AB     & GC, array, \& number work    & EhBASIC \\
\$AC--\$B9     & FAC1 \& FAC2 accumulators    & EhBASIC \\
\$BA--\$D7     & Byte-fetch subroutines       & EhBASIC \\
\$D8--\$DB     & PRNG state                   & EhBASIC \\
\$DC--\$E1     & NMI/IRQ handlers             & EhBASIC \\
\$E2           & HELP scratch                 & EhBASIC \\
\$E3--\$EE     & \textit{Free (12 bytes)}     & User \\
\$EF--\$FF     & Decimal conversion buffer    & EhBASIC \\
\bottomrule
\end{tabular}
\end{center}
